\section{Results}
\label{sec:results}

% Numbers come from `dockprepbench report` (paper/tables). Follow the benchmark-reporting
% conventions: denominators, both metrics, Wilson intervals, effect size before p-value,
% Holm-adjusted p. All four engines are complete.

\subsection{Baseline performance}

Of the 428 complexes, 412 were prepared and docked with Vina. The 16 failures all occurred
when converting the receptor to PDBQT: in 15, Meeko could not build a residue within
\SI{8}{\angstrom} of the box from its template, and in one it could not create the receptor
at all. These
complexes are counted as failures; the other engines do not use this step and docked all
428. With the baseline protocol, Vina placed the top-ranked pose within \SI{2}{\angstrom} of
the crystal ligand for \val{summary/vina/baseline/top1}\% [\val{summary/vina/baseline/top1_low}, \val{summary/vina/baseline/top1_high}] of complexes, and \val{summary/vina/baseline/pb_valid}\% [\val{summary/vina/baseline/pb_valid_low}, \val{summary/vina/baseline/pb_valid_high}] were also
\pbvalid{}.

The PoseBusters authors also docked these 428 complexes with Vina and released their
per-complex results with the dataset \citep{posebusters2023data}. Their run reached \val{posebusters_published/all/theirs_top1}\%
[\val{posebusters_published/all/theirs_top1_low}, \val{posebusters_published/all/theirs_top1_high}] top-1 success and \val{posebusters_published/all/theirs_pb_valid}\% \pbvalid{} success, 6.1 points above our baseline
(their run succeeded on 49 complexes where ours failed, ours on 23 where theirs failed;
exact McNemar $p = 0.003$). The two runs agreed on 356 of the 428 complexes. Our
preparation failures account for little of the difference: their run solved 5 of those 16
complexes, and on the 412 we docked the gap was 5.1 points (\val{posebusters_published/docked/theirs_top1}\% against \val{posebusters_published/docked/ours_top1}\%,
$p = \val{posebusters_published/docked/p}$). The protocols differ in several of the choices studied here. Theirs used Vina
1.2.3 with exhaustiveness 32 and 40 poses, a cubic search box of side \SI{25}{\angstrom}
centred on the crystal ligand, and receptors protonated with Reduce \citep{word1999reduce}
and converted with the prepare\_receptor script of ADFR \citep{ravindranath2015autodockfr}
(\citealp{buttenschoen2024posebusters}, ESI section S1), against exhaustiveness 8, nine
poses, a box
padded by \SI{4}{\angstrom} around the ligand and PDBFixer/OpenMM preparation here. Our
\SI{2}{\angstrom}-padding arm reached \val{comparisons/vina/top1/box_padding_2/arm_rate}\%, not significantly different from their result
($p = 0.82$). A difference of six points between two careful Vina protocols on the same
complexes is itself comparable to the preparation effects reported below.

At the baseline padding, scoring and search limited success about equally. A pose within
\SI{2}{\angstrom} was among the nine returned for 303 complexes (70.8\%), but ranked first
for only 198. Of the 230 complexes without a top-1 success, 105 were scoring failures (a
correct pose returned but not ranked first), while 109 had no correct pose among those
returned: 89 in the right site in the wrong binding mode and 20 in the wrong site; 16 failed
preparation; a further 19 top poses were correct but not \pbvalid{}.

\subsection{Run-to-run variation}

Repeating the Vina baseline with four docking seeds gave top-1 success of \val{summary/vina/baseline/top1}\%, \val{summary/vina/baseline_seed1/top1}\%,
\val{summary/vina/baseline_seed2/top1}\% and \val{summary/vina/baseline_seed3/top1}\% (mean \val{noise/vina/top1_mean}\%, standard deviation 0.4 points, range 0.9 points). The
aggregate rate is therefore stable to within one point. Individual outcomes are less stable:
two seeds disagree on \val{noise/vina/flips_mean} complexes on average (16--19 across the six pairs). Over the
four replicates, 183 complexes always succeeded, 212 always failed and 33 (7.7\%) changed
outcome with the seed alone. We use \val{noise/vina/flips_mean} changed complexes as a reference level for how
many outcomes an arm changes; it is a yardstick, not a significance threshold.

\subsection{Effect of each preparation choice}

\begin{figure}[t]
  \centering
  \includegraphics[width=0.92\textwidth]{engine_arms.pdf}
  \caption{Top-1 success (heavy-atom RMSD below \SI{2}{\angstrom}) for each preparation
    arm and engine on all 428 complexes, with 95\% Wilson intervals; preparation failures
    count as failures. Faint vertical lines mark each engine's baseline success. Drawn by
    \texttt{dockprepbench report --figures}.}
  \label{fig:success}
\end{figure}

Table~\ref{tab:comparisons} compares each arm with the Vina baseline on all 428 complexes.
Only one choice changed success beyond what multiple testing allows. A box padding of
\SI{2}{\angstrom} raised top-1 success from \val{comparisons/vina/top1/box_padding_2/baseline}\% [\val{summary/vina/baseline/top1_low}, \val{summary/vina/baseline/top1_high}] to \val{summary/vina/box_padding_2/top1}\% [\val{summary/vina/box_padding_2/top1_low}, \val{summary/vina/box_padding_2/top1_high}]
(\val{comparisons/vina/top1/box_padding_2/delta_pts} points [\val{comparisons/vina/top1/box_padding_2/delta_low}, \val{comparisons/vina/top1/box_padding_2/delta_high}]; 41 complexes gained, 12 lost; Holm $p = \val{comparisons/vina/top1/box_padding_2/p_holm}$) and \pbvalid{}
success from \val{comparisons/vina/pb_valid/box_padding_2/baseline}\% to \val{comparisons/vina/pb_valid/box_padding_2/arm_rate}\% (\val{comparisons/vina/pb_valid/box_padding_2/delta_pts} points; 41 gained, 15 lost; Holm $p = 0.003$). Larger
boxes lowered success slightly: \SI{6}{\angstrom} padding to \val{summary/vina/box_padding_6/top1}\% (\val{comparisons/vina/top1/box_padding_6/delta_pts} points; 10
gained, 18 lost; Holm $p = \val{comparisons/vina/top1/box_padding_6/p_holm}$) and \SI{8}{\angstrom} to \val{summary/vina/box_padding_8/top1}\% (\val{comparisons/vina/top1/box_padding_8/delta_pts} points; 19
gained, 28 lost; Holm $p = \val{comparisons/vina/top1/box_padding_6/p_holm}$).

The ligand choices moved success upwards without reaching significance after correction.
Assigning the dominant protonation state at pH 7.4 gave \val{summary/vina/ligand_dimorphite/top1}\% (\val{comparisons/vina/top1/ligand_dimorphite/delta_pts} points; 26 gained, 14
lost; unadjusted $p = 0.08$, Holm $p = \val{comparisons/vina/top1/ligand_dimorphite/p_holm}$); its \pbvalid{} effect was slightly larger
(\val{comparisons/vina/pb_valid/ligand_dimorphite/delta_pts} points; unadjusted $p = 0.04$, Holm $p = \val{comparisons/vina/pb_valid/ligand_dimorphite/p_holm}$). Regenerating the start conformer with
RDKit gave \val{summary/vina/ligand_rdkit_conformer/top1}\% (\val{comparisons/vina/top1/ligand_rdkit_conformer/delta_pts} points; 52 gained, 33 lost; unadjusted $p = 0.05$, Holm
$p = \val{comparisons/vina/top1/ligand_rdkit_conformer/p_holm}$).

\begin{table}[h]
  \centering
  \caption{Paired comparison of each arm with the Vina baseline on all 428 complexes
    (failures counted as failures). Top-1 success with 95\% Wilson interval; $\Delta$ with
    a 95\% interval for the paired difference (Newcombe); gained and lost
    are complexes that succeed in only the arm or only the baseline; $p$ is the exact
    McNemar test, Holm-adjusted over the five arms. Changed: complexes whose outcome
    differs from the baseline, and the ratio to the \val{noise/vina/flips_mean} changed by a change of seed.}
  \label{tab:comparisons}
  \small
  \setlength{\tabcolsep}{4pt}
  \begin{tabular}{lrrrrrr}
    \toprule
    Arm & Top-1 (\%) & $\Delta$ (pts) & Gained & Lost & Holm $p$ & Changed \\
    \midrule
    Baseline & 46.3 [41.6, 51.0] & -- & -- & -- & -- & -- \\
    Padding \SI{2}{\angstrom} & 53.0 [48.3, 57.7] & +6.8 [3.5, 10.0] & 41 & 12 & $4.1\times10^{-4}$
      & 53 (3.0$\times$) \\
    Padding \SI{6}{\angstrom} & 44.4 [39.8, 49.1] & $-1.9$ [$-4.3$, 0.5] & 10 & 18 & 0.37 & 28 (1.6$\times$) \\
    Padding \SI{8}{\angstrom} & 44.2 [39.5, 48.9] & $-2.1$ [$-5.2$, 1.0] & 19 & 28 & 0.37 & 47 (2.7$\times$) \\
    Protonation (pH 7.4) & 49.1 [44.4, 53.8] & +2.8 [$-0.1$, 5.7] & 26 & 14 & 0.24 & 40 (2.3$\times$) \\
    RDKit conformer & 50.7 [46.0, 55.4] & +4.4 [0.2, 8.6] & 52 & 33 & 0.20 & 85 (4.9$\times$) \\
    \bottomrule
  \end{tabular}
\end{table}

Every Vina arm changed the outcome of more complexes than a change of seed did
(Table~\ref{tab:comparisons}, last column), though for the \SI{6}{\angstrom} box only just
(28 against 16--19 between pairs of seeds). The clearest case is the start conformer: it
changed 85 outcomes, 4.9 times the seed level, although its net effect was only \val{comparisons/vina/top1/ligand_rdkit_conformer/delta_pts}
points. These ratios are descriptive: they compare counts of changed outcomes and do not
test whether one source of change matters more than the other. Most of these changes cancel, so a small net difference hides a large reshuffling of
which complexes succeed. Net success rates alone therefore understate how much a preparation
choice matters for any individual complex.

\subsection{How outcomes change: failure modes}

The failure classes below describe where top poses lay, not why they were chosen. By that
description, the tight box acted mainly through search. Moving from \SI{4}{\angstrom} to
\SI{2}{\angstrom} padding raised the number of complexes with any pose below
\SI{2}{\angstrom} from 303 to 338, and wrong-site top poses fell from 20 to 3. The largest
movements were wrong binding modes that became scoring failures (30 complexes, now with a
good pose that was not ranked first), scoring failures that became successes (28) and
wrong-site poses that moved into the site (17, 12 of them into a wrong binding mode). Larger boxes did the reverse: wrong-site top
poses rose to 45 at \SI{6}{\angstrom} and 64 at \SI{8}{\angstrom}, and complexes with any
good pose fell to 282 and 278. At \SI{8}{\angstrom}, 23 wrong-mode and 19 scoring failures
became wrong-site failures. Within the same scoring function, enlarging the box lets the
search find high-scoring poses outside the true site.

The RDKit start conformer acted differently. Its gains came from scoring failures (27) and
wrong binding modes (18) that became successes, but these were offset by successes that
became scoring failures (19) or wrong modes (11), and by scoring failures that became wrong
modes (19). The start conformer thus changes which basin Vina's search ends in, in both
directions, consistent with a sampling effect rather than a systematic improvement.
Ligand protonation mainly turned wrong binding modes into scoring failures (21) and scoring
failures into successes (15); eight correct but invalid poses became \pbvalid{}.

\subsection{Preparation effects across engines}

Table~\ref{tab:engines} gives the same comparisons for Smina, Gnina and rDock. Smina, which reproduces Vina's scoring function, behaved like Vina
throughout. Its baseline success was \val{summary/smina/baseline/top1}\% [\val{summary/smina/baseline/top1_low}, \val{summary/smina/baseline/top1_high}], and on all 428 complexes it did
not differ from Vina's (17 complexes gained, 14 lost; exact McNemar $p = \val{engines/smina/vina/all/p}$). The
\SI{2}{\angstrom} box was again the only significant arm, raising success to \val{summary/smina/box_padding_2/top1}\%
[\val{summary/smina/box_padding_2/top1_low}, \val{summary/smina/box_padding_2/top1_high}] (\val{comparisons/smina/top1/box_padding_2/delta_pts} points [\val{comparisons/smina/top1/box_padding_2/delta_low}, \val{comparisons/smina/top1/box_padding_2/delta_high}]; 43 gained, 10 lost; Holm $p = \val{comparisons/smina/top1/box_padding_2/p_holm}$), and larger
boxes again lowered it slightly (\SI{6}{\angstrom}: \val{comparisons/smina/top1/box_padding_6/delta_pts} points; \SI{8}{\angstrom}:
\val{comparisons/smina/top1/box_padding_8/delta_pts} points; Holm $p \ge \val{comparisons/smina/top1/box_padding_8/p_holm}$).

Gnina, which searches with the same machinery but re-ranks poses with a convolutional
neural network, reached \val{summary/gnina/baseline/top1}\% [\val{summary/gnina/baseline/top1_low}, \val{summary/gnina/baseline/top1_high}] at baseline, 16.4 points above Vina (103
complexes gained, 33 lost; $p = \val{engines/gnina/vina/all/p}$). The difference is not an artefact of
Vina's 16 preparation failures: on the 412 complexes both engines docked, success was
\val{engines/gnina/vina/shared/rate_b}\% for Vina and \val{engines/gnina/vina/shared/rate_a}\% for Gnina ($p = \val{engines/gnina/vina/shared/p}$). No preparation choice
changed Gnina's success significantly. The \SI{2}{\angstrom} box gave \val{comparisons/gnina/top1/box_padding_2/delta_pts} points [$-0.5$, \val{comparisons/gnina/top1/box_padding_2/delta_high}] and the
\SI{8}{\angstrom} box \val{comparisons/gnina/top1/box_padding_8/delta_pts} points (Holm $p = \val{comparisons/gnina/top1/box_padding_2/p_holm}$ and \val{comparisons/gnina/top1/box_padding_8/p_holm}); ligand protonation gave
\val{comparisons/gnina/top1/ligand_dimorphite/delta_pts} points and the RDKit conformer \val{comparisons/gnina/top1/ligand_rdkit_conformer/delta_pts} points.

rDock, whose search and scoring share nothing with the Vina family, reached \val{summary/rdock/baseline/top1}\% [\val{summary/rdock/baseline/top1_low},
\val{summary/rdock/baseline/top1_high}] at baseline, 16.8 points below Vina (33 complexes gained, 105 lost;
$p = \val{engines/rdock/vina/all/p}$) and 33.2 points below Gnina. It was by far the most sensitive to
the search space. The \SI{2}{\angstrom} box raised its success to \val{comparisons/rdock/top1/box_padding_2/arm_rate}\% [\val{summary/rdock/box_padding_2/top1_low}, \val{summary/rdock/box_padding_2/top1_high}]
(\val{comparisons/rdock/top1/box_padding_2/delta_pts} points [\val{comparisons/rdock/top1/box_padding_2/delta_low}, \val{comparisons/rdock/top1/box_padding_2/delta_high}]; 62 gained, 12 lost; Holm $p = \val{comparisons/rdock/top1/box_padding_2/p_holm}$), and larger boxes lowered
it to \val{comparisons/rdock/top1/box_padding_6/arm_rate}\% at \SI{6}{\angstrom} (\val{comparisons/rdock/top1/box_padding_6/delta_pts} points; 7 gained, 48 lost; Holm
$p = \val{comparisons/rdock/top1/box_padding_6/p_holm}$) and \val{comparisons/rdock/top1/box_padding_8/arm_rate}\% at \SI{8}{\angstrom} (\val{comparisons/rdock/top1/box_padding_8/delta_pts} points; 7 gained, 59
lost; Holm $p = \val{comparisons/rdock/top1/box_padding_8/p_holm}$). The ligand choices did not change its success
significantly: protonation gave +0.2 points and the RDKit conformer \val{comparisons/rdock/top1/ligand_rdkit_conformer/delta_pts} points (Holm
$p = 1$ and \val{comparisons/rdock/top1/ligand_rdkit_conformer/p_holm}).

The geometric pattern behind the box effect was shared. Every engine placed more top poses in
the wrong site as the box grew: from 20, 30, 4 and 91 complexes at \SI{4}{\angstrom}
padding (Vina, Smina, Gnina, rDock) to 64, 72, 36 and 219 at \SI{8}{\angstrom}, and down
to 3, 4, 2 and 8 at \SI{2}{\angstrom}. Gnina's re-scoring therefore removes most, but not
all, of the wrong-site errors that a larger box invites, while rDock, whose cavity grows
with the box, is the most exposed to them: the tight box raised the number of complexes for
which any of its poses lay within \SI{2}{\angstrom} of the crystal ligand from 186 to 271.

The spread of top-1 success over the six protocols (baseline and five arms) was 8.9 points for Vina
(44.2--\val{summary/vina/box_padding_2/top1}\%), 10.5 points for Smina (44.2--\val{summary/smina/box_padding_2/top1}\%), 7.0 points for Gnina
(\val{summary/gnina/box_padding_8/top1}--65.7\%) and 23.8 points for rDock (\val{summary/rdock/box_padding_8/top1}--41.1\%). At the baseline protocol the
engines spanned 33.2 points, from rDock to Gnina, and Gnina led Vina by 16.4. These spreads
run from the best to the worst protocol; they are not the effect of any one choice, which is
the paired difference reported above. The engine differences are differences between complete
pipelines, each with its own receptor representation, search and scoring
(Table~\ref{tab:receptor-paths}). For the Vina family, choosing the re-scoring pipeline
thus mattered about twice as much as any single preparation choice, yet the best preparation arm closed about 40\% of the gap between Vina
and Gnina (6.8 of 16.4 points). For rDock, the box alone moved success by more than its
16.8-point gap to Vina.

\begin{table}[h]
  \centering
  \caption{Top-1 success by engine and arm on all 428 complexes, with 95\% Wilson
    intervals; $\Delta$ is the change from that engine's baseline and $p$ is the exact
    McNemar test, Holm-adjusted over the five arms of each engine.}
  \label{tab:engines}
  \footnotesize
  \setlength{\tabcolsep}{3.5pt}
  \resizebox{\textwidth}{!}{%
  \begin{tabular}{lrrrrrr}
    \toprule
    & \multicolumn{2}{c}{Smina} & \multicolumn{2}{c}{Gnina} & \multicolumn{2}{c}{rDock} \\
    \cmidrule(lr){2-3} \cmidrule(lr){4-5} \cmidrule(lr){6-7}
    Arm & Top-1 (\%) & $\Delta$ (Holm $p$) & Top-1 (\%) & $\Delta$ (Holm $p$) & Top-1 (\%) & $\Delta$ (Holm $p$) \\
    \midrule
    Baseline & 47.0 [42.3, 51.7] & -- & 62.6 [57.9, 67.1] & -- & 29.4 [25.3, 33.9] & -- \\
    Padding \SI{2}{\angstrom} & 54.7 [49.9, 59.3] & +7.7 ($2.8\times10^{-5}$) & 65.7 [61.0, 70.0] & +3.0 (0.47) & 41.1 [36.6, 45.8] & +11.7 ($1.1\times10^{-8}$) \\
    Padding \SI{6}{\angstrom} & 45.6 [40.9, 50.3] & $-1.4$ (0.69) & 60.3 [55.6, 64.8] & $-2.3$ (0.47) & 19.9 [16.4, 23.9] & $-9.6$ ($3.9\times10^{-8}$) \\
    Padding \SI{8}{\angstrom} & 44.2 [39.5, 48.9] & $-2.8$ (0.38) & 58.6 [53.9, 63.2] & $-4.0$ (0.30) & 17.3 [14.0, 21.2] & $-12.1$ ($1.2\times10^{-10}$) \\
    Protonation (pH 7.4) & 48.1 [43.4, 52.9] & +1.2 (0.69) & 59.3 [54.6, 63.9] & $-3.3$ (0.47) & 29.7 [25.5, 34.2] & +0.2 (1) \\
    RDKit conformer & 50.2 [45.5, 54.9] & +3.3 (0.45) & 63.3 [58.7, 67.7] & +0.7 (0.83) & 25.5 [21.6, 29.8] & $-4.0$ (0.086) \\
    \bottomrule
  \end{tabular}}
\end{table}

\subsection{Robustness on the 308-complex subset}

Restricting every comparison to the 308 complexes on which the journal version of
PoseBusters reports results (those whose ligand has no crystal contact, that is, no atom
within \SI{5}{\angstrom} of a symmetry-related copy of the protein)
raised the absolute rates of the Vina-family engines and left the top-1 conclusions
unchanged. Baseline top-1 success was
\val{subset_308/summary/vina/baseline/top1}\% [\val{subset_308/summary/vina/baseline/top1_low}, \val{subset_308/summary/vina/baseline/top1_high}] for Vina, \val{subset_308/summary/smina/baseline/top1}\% [\val{subset_308/summary/smina/baseline/top1_low}, \val{subset_308/summary/smina/baseline/top1_high}] for Smina and \val{subset_308/summary/gnina/baseline/top1}\% [\val{subset_308/summary/gnina/baseline/top1_low}, \val{subset_308/summary/gnina/baseline/top1_high}] for
Gnina, between 6 and 7 points higher than on all 428 complexes, so the excluded complexes
were harder for these engines. Gnina's lead over Vina was 15.9 points
($p = 1.4 \times 10^{-6}$). The \SI{2}{\angstrom} box remained the only significant arm
for Vina (\val{subset_308/comparisons/vina/top1/box_padding_2/delta_pts} points; 25 gained, 8 lost; Holm $p = \val{subset_308/comparisons/vina/top1/box_padding_2/p_holm}$) and Smina (\val{subset_308/comparisons/smina/top1/box_padding_2/delta_pts} points; Holm
$p = \val{subset_308/comparisons/smina/top1/box_padding_2/p_holm}$), and no arm was significant for Gnina. The one notable shift was Gnina's
protonation arm, which lowered success by 6.2 points on this subset (19 gained, 38 lost;
Holm $p = 0.08$) against 3.3 points on all 428. rDock's baseline was \val{subset_308/summary/rdock/baseline/top1}\% on the
subset, close to its rate on all 428, and its box effects were as large: \val{subset_308/comparisons/rdock/top1/box_padding_2/delta_pts} points with
the \SI{2}{\angstrom} box and \val{subset_308/comparisons/rdock/top1/box_padding_6/delta_pts} and \val{subset_308/comparisons/rdock/top1/box_padding_8/delta_pts} points at 6 and \SI{8}{\angstrom} (Holm
$p \le \val{subset_308/comparisons/rdock/top1/box_padding_6/p_holm}$), with neither ligand choice significant for top-1 success.
For \pbvalid{} success, the subset changed two conclusions: the \SI{2}{\angstrom} box was no
longer significant for Vina (Holm $p = \val{subset_308/comparisons/vina/pb_valid/box_padding_2/p_holm}$) or Smina (\val{subset_308/comparisons/smina/pb_valid/box_padding_2/p_holm}), and the RDKit conformer
lowered rDock's \pbvalid{} success significantly (\val{subset_308/comparisons/rdock/pb_valid/ligand_rdkit_conformer/delta_pts} points; Holm $p = \val{subset_308/comparisons/rdock/pb_valid/ligand_rdkit_conformer/p_holm}$).
This subset is a robustness check within the same dataset, not an independent validation.

\subsection{Robustness to the wrong-site threshold and to removed cofactors}
\label{sec:robustness}

The box-size pattern does not depend on the \SI{4}{\angstrom} centroid threshold that
separates wrong sites from wrong binding modes. At every threshold from 2 to
\SI{6}{\angstrom}, and for every engine, the \SI{2}{\angstrom} box placed fewer top poses in
the wrong site than the baseline and the \SI{8}{\angstrom} box placed more
(Table~\ref{tab:threshold}); only the share of misses called wrong site changes with the
threshold.

\begin{table}[h]
  \centering
  \caption{Top-1 poses classed as wrong site at three centroid thresholds, for the baseline
    (\SI{4}{\angstrom} padding) and the \SI{2}{\angstrom} and \SI{8}{\angstrom} boxes (428
    complexes). Counts at 3 and \SI{5}{\angstrom} follow the same pattern.}
  \label{tab:threshold}
  \small
  \begin{tabular}{lccc}
    \toprule
    Engine & \SI{2}{\angstrom} threshold & \SI{4}{\angstrom} threshold & \SI{6}{\angstrom} threshold \\
     & 2 / 4 / \SI{8}{\angstrom} box & 2 / 4 / \SI{8}{\angstrom} box & 2 / 4 / \SI{8}{\angstrom} box \\
    \midrule
    Vina & 27 / 59 / 92 & 3 / 20 / 64 & 0 / 4 / 41 \\
    Smina & 31 / 68 / 96 & 4 / 30 / 72 & 1 / 10 / 41 \\
    Gnina & 17 / 31 / 58 & 2 / 4 / 36 & 0 / 1 / 17 \\
    rDock & 68 / 169 / 257 & 8 / 91 / 219 & 0 / 30 / 192 \\
    \bottomrule
  \end{tabular}
\end{table}

Removing cofactors affects the PoseBusters checks more than it affects success. Of the 428
complexes, 244 contain a heterogen other than common additives (organic cofactors, metal
and other ions). Among them, a correct top pose (RMSD below \SI{2}{\angstrom}) failed
PoseBusters only on a cofactor check for 13 complexes with Vina, 12 with Smina, 24 with Gnina
and 7 with rDock, poses placed where a removed cofactor would be. On the 184 complexes
without such heterogens, baseline top-1 success was \val{robust_no_cofactor_summary/vina/baseline/top1}\% for Vina, \val{robust_no_cofactor_summary/smina/baseline/top1}\% for Smina,
\val{robust_no_cofactor_summary/gnina/baseline/top1}\% for Gnina and \val{robust_no_cofactor_summary/rdock/baseline/top1}\% for rDock, so the engines kept their order. The
\SI{2}{\angstrom} box raised success for all four (Vina \val{robust_no_cofactor_comparisons/vina/top1/box_padding_2/delta_pts} points [\val{robust_no_cofactor_comparisons/vina/top1/box_padding_2/delta_low}, \val{robust_no_cofactor_comparisons/vina/top1/box_padding_2/delta_high}]; Smina \val{robust_no_cofactor_comparisons/smina/top1/box_padding_2/delta_pts}
[\val{robust_no_cofactor_comparisons/smina/top1/box_padding_2/delta_low}, \val{robust_no_cofactor_comparisons/smina/top1/box_padding_2/delta_high}]; Gnina \val{robust_no_cofactor_comparisons/gnina/top1/box_padding_2/delta_pts} [\val{robust_no_cofactor_comparisons/gnina/top1/box_padding_2/delta_low}, 7.1]; rDock \val{robust_no_cofactor_comparisons/rdock/top1/box_padding_2/delta_pts} [\val{robust_no_cofactor_comparisons/rdock/top1/box_padding_2/delta_low}, \val{robust_no_cofactor_comparisons/rdock/top1/box_padding_2/delta_high}]), significantly after
correction only for rDock (Holm $p = \val{robust_no_cofactor_comparisons/rdock/top1/box_padding_2/p_holm}$) in this smaller set. One effect differed: the
RDKit start conformer helped more in complexes without cofactors (Vina \val{robust_no_cofactor_comparisons/vina/top1/ligand_rdkit_conformer/delta_pts} points
[\val{robust_no_cofactor_comparisons/vina/top1/ligand_rdkit_conformer/delta_low}, \val{robust_no_cofactor_comparisons/vina/top1/ligand_rdkit_conformer/delta_high}], Smina \val{robust_no_cofactor_comparisons/smina/top1/ligand_rdkit_conformer/delta_pts} [\val{robust_no_cofactor_comparisons/smina/top1/ligand_rdkit_conformer/delta_low}, \val{robust_no_cofactor_comparisons/smina/top1/ligand_rdkit_conformer/delta_high}], against +4.4 and \val{comparisons/smina/top1/ligand_rdkit_conformer/delta_pts} on all 428), an exploratory
observation that this design cannot explain.

\subsection{Which complexes are sensitive to preparation}

Part of every arm's effect falls on complexes that are unstable anyway. The 33 complexes
whose Vina outcome changed with the seed alone were strongly over-represented among each
arm's flips (odds ratios 10.5--37; Holm $p < 10^{-8}$ for every arm): with the
\SI{2}{\angstrom} box, 58\% of them flipped against 9\% of the other complexes. But they
account for a minority of most arms' flips (Table~\ref{tab:stable-flips}). The RDKit start
conformer changed the outcome of 63 complexes that never changed with the seed, the
\SI{2}{\angstrom} box 34 and protonation 24. These complexes did not change across four
seeds, so most of these flips are likely effects of the preparation choice itself, although
four replicates cannot rule out that some of them have a small chance of flipping by
chance.

\begin{table}[h]
  \centering
  \caption{Vina complexes whose top-1 outcome changed with each arm, split by whether
    the outcome also changed between seed replicates of the baseline (33 seed-sensitive
    complexes; 395 stable).}
  \label{tab:stable-flips}
  \begin{tabular}{lrrr}
    \toprule
    Arm & Flipped & of them seed-sensitive & of them stable \\
    \midrule
    Padding \SI{2}{\angstrom} & 53 & 19 & 34 \\
    Padding \SI{6}{\angstrom} & 28 & 17 & 11 \\
    Padding \SI{8}{\angstrom} & 47 & 19 & 28 \\
    Protonation (pH 7.4) & 40 & 16 & 24 \\
    RDKit conformer & 85 & 22 & 63 \\
    \bottomrule
  \end{tabular}
\end{table}

Ligand properties explained little. The one association that survived correction was with
the \SI{2}{\angstrom} box for Smina, whose flips concentrated on more flexible ligands: a
median of 10 rotatable bonds against 6 (rank-biserial $r = 0.31$, Holm $p = 0.001$). Vina
showed the same tendency without reaching significance (median 29 heavy atoms against 24,
$r = 0.21$, Holm $p = 0.057$). Protonation flips leaned towards ligands whose charge it
changed for the three Vina-family engines (Vina: 80\% of flipped against 62\% of unflipped
complexes, odds ratio 2.5, $p = 0.024$; Gnina 1.7; Smina 1.6) but not for rDock (0.87),
and no engine's association survived correction (Holm $p \ge 0.12$). Flips from the RDKit conformer showed no association with ligand size
or flexibility in any engine, and no arm's flips were associated with the presence of a
cofactor. All but 10 of the supplied ligands are neutral, so the formal charge of the
supplied ligand itself could not be tested meaningfully.

\subsection{Interaction between protonation and box size, and a box sized to the ligand}
\label{sec:interaction}

Because each arm changes one choice from the baseline, the effects above are measured at
the baseline setting of the others. Two further sets of Vina runs tested that assumption.
First, protonation at pH 7.4 was repeated at each box padding (Table~\ref{tab:interaction}).
Its effect declined steadily as the box grew, from \val{followup_interaction/2/diff_pts} points at \SI{2}{\angstrom} (30
complexes gained, 12 lost; Holm $p = \val{followup_interaction/2/p_holm}$ over the four paddings, a correction family
separate from the five primary arms) to \val{comparisons/vina/top1/ligand_dimorphite/delta_pts}, +1.4 and
$-0.5$ points at 4, 6 and \SI{8}{\angstrom}. Compared complex by complex with the effect at
the baseline padding, the change was not significant at any single padding (sign test on
the difference in effect, Holm $p \ge 0.46$), so the trend is consistent but not
established. It is what would be expected if a large box lets wrong-site poses dominate,
whatever the ligand's charge state, while in a tight box the charge state decides between
nearby poses. Combining the two best choices gave the best Vina protocol in this study:
\SI{2}{\angstrom} padding with protonation at pH 7.4 reached \val{followup_combined/box_padding_2+ligand_dimorphite/rate_b}\% [\val{followup_combined/box_padding_2+ligand_dimorphite/low_b}, \val{followup_combined/box_padding_2+ligand_dimorphite/high_b}], 11.0
points above the baseline (60 complexes gained, 13 lost; exact McNemar
$p = \val{followup_combined/box_padding_2+ligand_dimorphite/p}$) and two thirds of the way to Gnina's baseline (\val{summary/gnina/baseline/top1}\%).

\begin{table}[h]
  \centering
  \caption{Vina top-1 success with and without ligand protonation at pH 7.4, at each box
    padding (428 complexes; failures counted as failures; 95\% Wilson intervals). Gained
    and lost are relative to the same padding without protonation; Holm-adjusted exact
    McNemar $p$ over the four paddings.}
  \label{tab:interaction}
  \begin{tabular}{lrrrrr}
    \toprule
    Padding & Supplied ligand (\%) & Protonated (\%) & $\Delta$ (pts) & Gained / lost & Holm $p$ \\
    \midrule
    \SI{2}{\angstrom} & 53.0 [48.3, 57.7] & 57.2 [52.5, 61.8] & +4.2 & 30 / 12 & 0.032 \\
    \SI{4}{\angstrom} (baseline) & 46.3 [41.6, 51.0] & 49.1 [44.4, 53.8] & +2.8 & 26 / 14 & 0.24 \\
    \SI{6}{\angstrom} & 44.4 [39.8, 49.1] & 45.8 [41.1, 50.5] & +1.4 & 24 / 18 & 0.88 \\
    \SI{8}{\angstrom} & 44.2 [39.5, 48.9] & 43.7 [39.1, 48.4] & $-0.5$ & 27 / 29 & 0.89 \\
    \bottomrule
  \end{tabular}
\end{table}

Second, the box was sized to the ligand rather than padded around it: a cube 2.857 times
the radius of gyration of the docked conformer, centred on the crystal ligand, the
optimum \citet{feinstein2015boxsize} found for Vina. These boxes were smaller than the
baseline box for 427 complexes (median volume \val{followup_gyration_geometry/median_ratio_to_baseline} times) and smaller than the
\SI{2}{\angstrom} box for 331 (\val{followup_gyration_geometry/median_ratio_to_box2} times), and for 207 complexes part of the crystal
ligand lay outside the box. Success was \val{followup_gyration/vs_baseline/rate_b}\% [\val{followup_gyration/vs_baseline/low_b}, \val{followup_gyration/vs_box_padding_2/high_b}], above the baseline (\val{followup_gyration/vs_baseline/diff_pts}
points; 46 gained, 29 lost; exact McNemar $p = \val{followup_gyration/vs_baseline/p}$) but not above the
\SI{2}{\angstrom} box (\val{comparisons/vina/top1/box_padding_2/arm_rate}\%; $p = \val{followup_gyration/vs_box_padding_2/p}$). The result split cleanly on whether the
crystal ligand fitted (an exploratory, post hoc split): for the 221 complexes whose
crystal ligand lay wholly inside the box, success rose from \val{followup_gyration/vs_baseline_inside/rate_a}\% to \val{followup_gyration/vs_baseline_inside/rate_b}\% (22 gained,
4 lost), while for the 207 where it did not, success was unchanged (40.6\% and \val{followup_gyration/vs_baseline_outside/rate_b}\%; 24
gained, 25 lost). A box scaled to the ligand works well when it contains the binding
mode, and the gain is lost where the rule truncates it; a small fixed padding around the
reference ligand avoids that failure.

\subsection{Agreement across preparation arms as a candidate confidence signal}
\label{sec:agreement}

The sensitivity above can also be used. For each engine and complex, the baseline top-1
pose was compared with the top-1 pose of each of the five primary arms, and the number of
arms whose pose lay within \SI{2}{\angstrom} of it (0--5; a missing pose counts as
disagreement) was tested as a predictor of whether the baseline pose was correct
(Figure~\ref{fig:agreement}). When all five arms agreed, the baseline pose was correct for
\val{agreement_levels/vina/5/rate}\% [\val{agreement_levels/vina/5/low}, \val{agreement_levels/vina/5/high}] of 166 complexes with Vina, \val{agreement_levels/smina/5/rate}\% [\val{agreement_levels/smina/5/low}, \val{agreement_levels/smina/5/high}] of 171 with Smina,
\val{agreement_levels/gnina/5/rate}\% [\val{agreement_levels/gnina/5/low}, \val{agreement_levels/gnina/5/high}] of 194 with Gnina and \val{agreement_levels/rdock/5/rate}\% [\val{agreement_levels/rdock/5/low}, \val{agreement_levels/rdock/5/high}] of 45 with rDock. When none
agreed, it was correct for at most \val{agreement_levels/vina/0/rate}\% (Vina, 2 of 36).

\begin{figure}[t]
  \centering
  \includegraphics[width=0.7\textwidth]{pose_agreement.pdf}
  \caption{Baseline top-1 success against the number of the five primary arms whose top-1
    pose lies within \SI{2}{\angstrom} of the baseline top-1 pose, for each engine (428
    complexes; failures counted as failures and missing poses as disagreement; 95\% Wilson
    intervals). Drawn by \texttt{python -m dockprepbench.analysis.agreement}.}
  \label{fig:agreement}
\end{figure}

As a ranking of poses by confidence, agreement across arms outperformed the docking score
for every engine, by 0.16 to 0.21 in the area under the ROC curve (AUROC;
Table~\ref{tab:agreement}), and it outperformed Gnina's own CNN pose score (\val{agreement_signals/prep_agree/gnina/auroc} against
\val{agreement_signals/cnn/gnina/auroc}; difference 0.051 [0.011, 0.093]). Agreement across engines, the established form of
consensus, was comparable: better for Smina (difference 0.048 [0.009, 0.087]), worse for
Gnina (\val{agreement_signals/engine_agree/gnina/diff_vs_prep} [\val{agreement_signals/engine_agree/gnina/diff_low_vs_prep}, \val{agreement_signals/engine_agree/gnina/diff_high_vs_prep}]), and not distinguishable for Vina and rDock. The two
carried partly different information: their sum was better than agreement across arms alone
for every engine (by 0.024 to 0.069, all intervals excluding zero), reaching 0.95 for
rDock.

The signal is not simply the stochasticity of the search. For Vina, agreement across the
three seed replicates of the baseline, whose preparation is identical, gave an AUROC of
\val{agreement_signals/seed_agree/vina/auroc} [\val{agreement_signals/seed_agree/vina/low}, \val{agreement_signals/seed_agree/vina/high}], not distinguishable from the docking score and 0.183 [0.148, 0.220]
below agreement across arms. A change of seed moved the top-1 pose by more than
\SI{2}{\angstrom} in \val{agreement_pose_change/vina/seeds_moved}\% of comparisons, a change of preparation in \val{agreement_pose_change/vina/arms_moved}\%: preparation
choices perturb the search more, and where the top pose survives that perturbation it is
usually right.

Most of the signal needs only two extra runs. Averaged over every pair of arms, agreement
gave an AUROC of 0.775 for Vina, 0.744 for Smina, \val{agreement_subsets/gnina/box_padding_2+ligand_rdkit_conformer/auroc_mean} for Gnina and 0.836 for rDock,
against \val{agreement_signals/prep_plus_engine/gnina/low}, \val{agreement_subsets/smina/box_padding_2+box_padding_6+box_padding_8+ligand_dimorphite+ligand_rdkit_conformer/auroc_mean}, \val{agreement_signals/prep_agree/gnina/auroc} and \val{agreement_subsets/rdock/box_padding_2+box_padding_6+box_padding_8+ligand_dimorphite+ligand_rdkit_conformer/auroc_mean} with all five arms; a single arm gave 0.67 to 0.76.
The same pair, the \SI{2}{\angstrom} box with the RDKit start conformer, was the best of the
ten for all four engines (\val{agreement_subsets/vina/box_padding_2+ligand_rdkit_conformer/auroc_max}, \val{agreement_subsets/smina/box_padding_2+ligand_rdkit_conformer/auroc_max}, \val{agreement_subsets/vina/box_padding_2+ligand_rdkit_conformer/auroc_max} and \val{agreement_subsets/rdock/box_padding_2+ligand_rdkit_conformer/auroc_max}), within 0.03 to 0.05 of all five
arms. It was selected on these data, so its advantage over the average pair is optimistic,
but that one pair came first for every engine suggests the choice is not arbitrary: it
combines a change of search space with a change of starting point.

Because the signal was found on these complexes, it was also evaluated out of sample
(Table~\ref{tab:agreement-cv}). Under cross-validation, the agreement levels learned on
training folds ranked held-out poses with an AUROC of \val{agreement_validation/vina/prep_agree/auroc} for Vina, \val{agreement_validation/smina/prep_agree/auroc} for Smina, \val{agreement_validation/gnina/prep_agree/auroc}
for Gnina and \val{agreement_validation/rdock/prep_agree/auroc} for rDock, against \val{agreement_validation/smina/score/auroc} to \val{agreement_validation/rdock/score/auroc} for the docking score, and with lower
Brier scores (\val{agreement_validation/rdock/prep_agree/brier} to \val{agreement_validation/gnina/best_pair/brier_hi} against \val{agreement_validation/rdock/score/brier} to \val{agreement_validation/smina/score/brier}). Agreement added to the score rather
than restating it: a model with both reached held-out AUROCs of \val{agreement_validation/smina/score+prep_agree/auroc} to \val{agreement_validation/rdock/score+prep_agree/auroc}, and adding
agreement to the score improved the fit for every engine (likelihood-ratio test, $p <
10^{-30}$). Choosing the best pair of arms inside each training fold, which removes the
optimism of choosing it on all the data, gave held-out AUROCs of \val{agreement_validation/smina/best_pair/auroc} to \val{agreement_validation/rdock/best_pair/auroc}, about 0.02
below the in-sample values above. Missing poses do not drive the signal: only 16 Vina and 2
rDock complexes had a missing arm pose, and without them Vina's AUROC was \val{agreement_missing/vina/prep_auroc_complete} instead of
\val{agreement_signals/prep_agree/vina/auroc}. All of this remains within one dataset; Section~\ref{sec:astex} tests it on another.

\begin{table}[h]
  \centering
  \caption{Held-out performance of confidence predictors for baseline top-1 success: mean
    over 20 repetitions of 5-fold cross-validation by complex (range over repetitions in
    brackets). Agreement: training-fold success rate at each agreement level. Best pair:
    the pair of arms with the highest training-fold AUROC. Lower Brier is better.}
  \label{tab:agreement-cv}
  \small
  \setlength{\tabcolsep}{4pt}
  \begin{tabular}{lrrrrr}
    \toprule
    Engine & Score & Agreement (5 arms) & Score + agreement & Best pair & Brier: score / agreement \\
    \midrule
    Vina & 0.671 & 0.838 [0.832, 0.845] & 0.872 & 0.809 [0.800, 0.818] & 0.225 / 0.147 \\
    Smina & 0.637 & 0.809 [0.801, 0.816] & 0.857 & 0.790 [0.781, 0.796] & 0.234 / 0.162 \\
    Gnina & 0.696 & 0.835 [0.826, 0.844] & 0.868 & 0.802 [0.784, 0.813] & 0.207 / 0.137 \\
    rDock & 0.717 & 0.884 [0.881, 0.887] & 0.906 & 0.871 [0.860, 0.883] & 0.201 / 0.118 \\
    \bottomrule
  \end{tabular}
\end{table}

\begin{table}[h]
  \centering
  \caption{AUROC for predicting whether the baseline top-1 pose is correct (RMSD below
    \SI{2}{\angstrom}), with 95\% bootstrap intervals (2,000 resamples of complexes; 428
    complexes per engine). Docking scores enter negated, so higher is more confident.
    Gnina's docking score is its minimised Vina-like affinity; its CNN score is shown
    separately. Agreement counts the seed replicates, primary arms or other engines whose
    top-1 pose lies within \SI{2}{\angstrom} of the baseline pose; seed replicates exist
    only for Vina.}
  \label{tab:agreement}
  \small
  \setlength{\tabcolsep}{4pt}
  \begin{tabular}{lrrrr}
    \toprule
    Signal & Vina & Smina & Gnina & rDock \\
    \midrule
    Docking score & 0.65 [0.60, 0.70] & 0.64 [0.59, 0.69] & 0.70 [0.65, 0.75] & 0.72 [0.67, 0.77] \\
    Gnina CNN score & -- & -- & 0.81 [0.77, 0.85] & -- \\
    Seeds agree (0--3) & 0.67 [0.64, 0.71] & -- & -- & -- \\
    Arms agree (0--5) & 0.86 [0.82, 0.89] & 0.83 [0.79, 0.87] & 0.86 [0.83, 0.90] & 0.90 [0.87, 0.93] \\
    Engines agree (0--3) & 0.89 [0.86, 0.92] & 0.88 [0.85, 0.91] & 0.78 [0.74, 0.82] & 0.92 [0.89, 0.95] \\
    Arms + engines & 0.91 [0.88, 0.94] & 0.90 [0.87, 0.93] & 0.89 [0.86, 0.92] & 0.95 [0.93, 0.97] \\
    \bottomrule
  \end{tabular}
\end{table}

Agreement can also choose the pose. Taking, among the six top-1 poses (baseline and five
arms), the one that most of the others agree with (ties keep the baseline pose) changed the
chosen pose for \val{agreement_consensus/top1/vina/changed_pose}\% (Vina) to \val{agreement_consensus/top1/gnina/changed_pose}\% (Gnina, rDock) of complexes and raised top-1 success
for every engine: Vina from \val{summary/vina/baseline/top1}\% to \val{agreement_consensus/top1/vina/consensus}\% (11 gained, 3 lost), Smina from \val{summary/smina/baseline/top1}\% to
\val{agreement_consensus/top1/smina/consensus}\% (8 gained, 1 lost), Gnina from \val{summary/gnina/baseline/top1}\% to \val{agreement_consensus/top1/gnina/consensus}\% (19 gained, 6 lost) and rDock from
\val{summary/rdock/baseline/top1}\% to \val{agreement_consensus/top1/rdock/consensus}\% (14 gained, 6 lost). PB-valid success moved the same way. The direction was
consistent, but no single engine's gain survived correction over the four engines (Holm
$p \ge 0.059$), so the consensus pose is a modest, unconfirmed improvement, whereas
agreement as a confidence signal is a large one. As tested, both used six runs per complex,
although the confidence signal keeps most of its value with three.

\subsection{Independent check on the Astex Diverse Set}
\label{sec:astex}

\result{Astex validation: H1 (2 \AA{} box), H2 (agreement over five arms vs score), H3 (the
pre-selected pair vs score) and H4 (success when all five arms agree), as specified in
docs/astex\_validation\_plan.md before the runs. To be filled in when the runs finish.}
